\documentclass{article}
\usepackage[utf8]{inputenc}
\usepackage{graphicx}
\usepackage{enumerate}
\usepackage{icomma}
\usepackage{amssymb}
\usepackage{amsmath}
\usepackage{tikz}
\usepackage{href-ul}
\hypersetup{
	colorlinks=true,
	linkcolor=blue,
	urlcolor=blue}
\usepackage{geometry}
\geometry{a4paper, top=15mm, left=15mm, right=15mm, bottom=15mm,
headsep=10mm, footskip=12mm}

\title{m11ablanfL246}

\begin{document}
\begin{tikzpicture}(10,3)
	\draw[ultra thick](2,0) --(10,0) -- (10,1.3) --(2,1.3) -- (2,0);
	\draw[fill=black](2,0)-- (0,.6) -- (2,.6) -- (2,0);
	\node at (6,.6) {\large Lösungsblatt von \href{https://www.okuyakl.de}{www.okuyakl.de}};
\end{tikzpicture}
\vspace{0.5 cm}

\noindent {\bf Aufgabe 1. a)}\\
Die Funktion ist an der Stelle, an der der Ausdruck zwischen den Betragsstrichen Null wird, also bei $x=1$ eventuell nicht differenzierbar.

\noindent {\bf Aufgabe 1. b)}\\
Es gilt: 
$$f(x)= \begin{cases} {1 \over (x-1) +1 } = {1\over x} & x \ge 1 \\
                      {1 \over -(x-1) + 1 } = { 1 \over 2 -x } & x <1
 \end{cases}
 $$
 $$
 \renewcommand{\arraystretch}{2}
 \begin{array}{lrcl}
 \underline{Fall 1.\quad x \ge 1}\\
 & \lim\limits_{x \to 1}\frac{{1\over x} - {1 \over 1}}{x-1}&=&
 \lim\limits_{x \to 1}\frac{{1-x \over x}}{x-1} \\
 & =\lim\limits_{x \to 1}-{1\over x}&=& -1   
 \end{array}
 $$
$$
\renewcommand{\arraystretch}{2}
 \begin{array}{lrcl}
 \underline{Fall 2.\quad x < 1}\\
 &\lim\limits_{x \to 1}\frac{{1\over 2-x} - {1 \over 2-1}}{x-1}&=&
 \lim\limits_{x \to 1}\frac{{1-2+x \over 2-x}}{x-1} \\
 &=\lim\limits_{x \to 1}{1\over 2-x}&=& 1   
 \end{array}
$$
Der linksseitige Wert des Differenzenquotienten stimmt nicht mit dem rechtsseitigen überein. Daraus folgt, dass die Funktion an der Stelle $x=1$ nicht differenzierbar ist.
 
\noindent {\bf Aufgabe 1. c)}\\
Die Funktion $f(x)=|x+3|$ ist an der Stelle $x=-3$ nicht differenzierbar.

\noindent {\bf Aufgabe 2. a)}\\
Die Steigung ist Null in den Punkten, wo die Funktion waagrecht verläuft, also 
$$f'(x)=0 \quad \textnormal{für} \quad x \in [4;6] \lor x = 8 $$

\noindent {\bf Aufgabe 2. b)}\\
Die Funktion ist nicht differenzierbar, wo sie Knickstellen aufweist, also in
$$x \in \{ 4; 6 \}$$

\noindent {\bf Aufgabe 3. a)}\\
\begin{minipage}{0.5\textwidth}
\includegraphics[width=6 cm]{reitan246}
\end{minipage}
\begin{minipage}{0.5\textwidth}
{\bf Aufgabe 3. b)}\\
$$
\begin{array}{rcl}
f'(x)&=& \lim\limits_{h \to 0}\frac{0,25(x+h)^2 - 0,25x^2}{h}\\
     &=& \lim\limits_{h \to 0}\frac{0,25x^2 + 0,5xh + 0,25 h^2 - 0,25x^2}{h}\\
     &=& \lim\limits_{h \to 0}\frac{h(0,5x +h)}{h}\\
     &=& \lim\limits_{h \to 0} 0,5x + h = 0,5x\\
\end{array}
$$
\end{minipage}

\begin{minipage}{0.5\textwidth}
\noindent {\bf Aufgabe 3. c)}\\
$$
\begin{array}{|rcl|rcl|}
\hline
f'(-3)&=& -1,5 & f'(-1) &=& -0,5 \\
f'(0) &=& 0  & f'(1) &=& 0,5 \\
f'(3)&=& 1,5  & & & \\
\hline 
\end{array}
$$
\end{minipage}
\begin{minipage}{0.5\textwidth}
\noindent {\bf Aufgabe 3. d,e)}\\
\includegraphics[width=6 cm]{abrei246}
\end{minipage}

\begin{minipage}{0.5\textwidth}
\noindent {\bf Aufgabe 4. b)}\\
$$
f'(x)= \begin{cases} 1 & x \ge 0 \\
                    -1 & x < 0 \end{cases}
$$
\end{minipage}
\begin{minipage}{0.5\textwidth}
\noindent {\bf Aufgabe 4. a)}\\
\includegraphics[width=6 cm]{absab246}
\end{minipage}

\begin{center}
	\includegraphics[width=7 cm]{../../viecher/endcomic.pdf}
	
	Hier geht es zurück zum \href{https://www.okuyakl.de/math/m11ablanfL246/aa246.pdf}{Aufgabenblatt}
\end{center}

%-Ende
\end{document}
